
\documentclass{article}
\usepackage{colortbl}
\usepackage{makecell}
\usepackage{multirow}
\usepackage{supertabular}

\begin{document}

\newcounter{utterance}

\centering \large Interaction Transcript for game `referencegame', experiment `random\_grids', episode 11 with reflexion{-}llama{-}3.1{-}8b{-}instruct.
\vspace{24pt}

{ \footnotesize  \setcounter{utterance}{1}
\setlength{\tabcolsep}{0pt}
\begin{supertabular}{c@{$\;$}|p{.15\linewidth}@{}p{.15\linewidth}p{.15\linewidth}p{.15\linewidth}p{.15\linewidth}p{.15\linewidth}}
    \# & $\;$A & \multicolumn{4}{c}{Game Master} & $\;\:$B\\
    \hline

    \theutterance \stepcounter{utterance}  
    & & \multicolumn{4}{p{0.6\linewidth}}{
        \cellcolor[rgb]{0.9,0.9,0.9}{
            \makecell[{{p{\linewidth}}}]{
                \texttt{\tiny{[P1$\langle$GM]}}
                \texttt{You have attempted this same task before and did not succeed. These are the plans you wrote after those attempts.} \\
\texttt{Trial \#1: New plan: I failed to distinguish the target grid from the distractors because my referring expression was too specific and focused on a single row. What I should have done is to look for a unique feature that is present in the target grid but not in the distractors. In this case, I should have noticed that the target grid has a unique "X" in the top{-}left corner, which is not present in the distractors. Next time, I will look for such unique features and incorporate them into my referring expression.} \\
\\ 
\texttt{Trial \#2: New plan: I failed to distinguish the target grid from the distractors because my referring expression was too specific and focused on a single row. What I should have done is to look for a unique feature that is present in the target grid but not in the distractors. In this case, I should have noticed that the target grid has a unique "X" in the top{-}left corner, which is not present in the distractors. However, I also failed to notice that the target grid has a unique combination of "X"s in the second row from the top, which is not present in the distractors. Next time, I will look for such unique features and incorporate them into my referring expression. I will also make sure to consider multiple features and their combinations to increase the chances of distinguishing the target grid from the distractors.} \\
\\ 
\\ 
\texttt{You are given three grids, where each of them is 5 by 5 in size.} \\
\texttt{Grids have empty cells marked with "▢" and filled cells marked with "X".} \\
\texttt{Your task is to generate a referring expression that best describes the target grid while distinguishing it from the two other distractor grids.} \\
\texttt{The first grid is the target grid, and the following two grids are the distractors.} \\
\\ 
\texttt{Target grid:} \\
\\ 
\texttt{X □ X X □} \\
\texttt{□ X □ X X} \\
\texttt{X X □ □ X} \\
\texttt{□ X X □ X} \\
\texttt{X □ X □ □} \\
\\ 
\texttt{Distractor grid 1:} \\
\\ 
\texttt{□ □ X X □} \\
\texttt{X □ □ X X} \\
\texttt{□ X X □ X} \\
\texttt{X X □ □ X} \\
\texttt{□ X X □ □} \\
\\ 
\texttt{Distractor grid 2:} \\
\\ 
\texttt{X X □ X □} \\
\texttt{□ X X □ X} \\
\texttt{X □ □ X □} \\
\texttt{□ X □ □ X} \\
\texttt{X □ X □ X} \\
\\ 
\texttt{Instruction: Describe the target grid.} \\
\texttt{Generate the referring expression starting with the tag "Expression: " for the given target grid. Omit any other text.} \\
            }
        }
    }
    & \\ \\

    \theutterance \stepcounter{utterance}  
    & \multicolumn{4}{p{0.6\linewidth}}{
        \cellcolor[rgb]{0.9,0.9,0.9}{
            \makecell[{{p{\linewidth}}}]{
                \texttt{\tiny{[P1$\rangle$GM]}}
                \texttt{Expression: The grid that has an "X" in the top{-}left corner and a unique combination of "X"s in the second row from the top, where the second row from the top has an "X" in the second column and an "X" in the fifth column.} \\
            }
        }
    }
    & & \\ \\

    \theutterance \stepcounter{utterance}  
    & & & \multicolumn{2}{p{0.3\linewidth}}{
        \cellcolor[rgb]{0.9,0.9,0.9}{
            \makecell[{{p{\linewidth}}}]{
                \texttt{\tiny{[GM$|$GM]}}
                \texttt{Expression: The grid that has an "X" in the top{-}left corner and a unique combination of "X"s in the second row from the top, where the second row from the top has an "X" in the second column and an "X" in the fifth column.} \\
            }
        }
    }
    & & \\ \\

    \theutterance \stepcounter{utterance}  
    & & \multicolumn{4}{p{0.6\linewidth}}{
        \cellcolor[rgb]{0.9,0.9,0.9}{
            \makecell[{{p{\linewidth}}}]{
                \texttt{\tiny{[GM$\rangle$P2]}}
                \texttt{You have attempted this same task before and did not succeed. These are the plans you wrote after those attempts.} \\
\texttt{Trial \#1: New plan: I failed to match the grid with the given expression because I didn't carefully examine the first row of each grid. Next time, I will pay closer attention to the first row of each grid and select the grid that matches the expression. In this case, the expression states that the grid has exactly two X's in the first row. I will re{-}examine the grids and select the one that meets this condition.} \\
\\ 
\texttt{Trial \#2: New plan: I failed to match the grid with the given expression because I didn't carefully examine the third row of each grid. The expression states that the grid has a row with two "X"s in the third row from the top, but I didn't check this condition thoroughly. Next time, I will pay closer attention to the third row of each grid and select the grid that matches the expression. I will re{-}examine the grids and select the one that meets this condition.} \\
\\ 
\\ 
\texttt{You are given three grids, where each of them is 5 by 5 in size.} \\
\texttt{Grids have empty cells marked with "▢" and filled cells marked with "X".} \\
\texttt{You are also given a referring expression that describes one of the given grids.} \\
\texttt{Your task is to select the grid that matches the given referring expression.} \\
\texttt{Generate only the number (in text) of the grid that matches the given expression by selecting first, second, or third.} \\
\\ 
\texttt{First:} \\
\\ 
\texttt{X X □ X □} \\
\texttt{□ X X □ X} \\
\texttt{X □ □ X □} \\
\texttt{□ X □ □ X} \\
\texttt{X □ X □ X} \\
\\ 
\texttt{Second:} \\
\\ 
\texttt{□ □ X X □} \\
\texttt{X □ □ X X} \\
\texttt{□ X X □ X} \\
\texttt{X X □ □ X} \\
\texttt{□ X X □ □} \\
\\ 
\texttt{Third:} \\
\\ 
\texttt{X □ X X □} \\
\texttt{□ X □ X X} \\
\texttt{X X □ □ X} \\
\texttt{□ X X □ X} \\
\texttt{X □ X □ □} \\
\\ 
\texttt{Expression: The grid that has an "X" in the top{-}left corner and a unique combination of "X"s in the second row from the top, where the second row from the top has an "X" in the second column and an "X" in the fifth column.} \\
\texttt{Question: Which grid does the expression refer to?} \\
\texttt{Start with the tag "Answer: ", followed by your selection. Omit any other text.} \\
            }
        }
    }
    & \\ \\

    \theutterance \stepcounter{utterance}  
    & & & \multicolumn{4}{p{0.6\linewidth}}{
        \cellcolor[rgb]{0.9,0.9,0.9}{
            \makecell[{{p{\linewidth}}}]{
                \texttt{\tiny{[GM$\langle$P2]}}
                \texttt{Answer: 1} \\
            }
        }
    }
     \\ \\

    \theutterance \stepcounter{utterance}  
    & & & \multicolumn{2}{p{0.3\linewidth}}{
        \cellcolor[rgb]{0.9,0.9,0.9}{
            \makecell[{{p{\linewidth}}}]{
                \texttt{\tiny{[GM$|$GM]}}
                \texttt{1} \\
            }
        }
    }
    & & \\ \\

    \theutterance \stepcounter{utterance}  
    & \multicolumn{4}{p{0.6\linewidth}}{
        \cellcolor[rgb]{0.95,0.95,0.95}{
            \makecell[{{p{\linewidth}}}]{
                \texttt{\tiny{[P2$\rightarrow$P1]}}
                \texttt{Post{-}episode Reflexion (not sent to GM):} \\
\texttt{New plan: I failed to distinguish the target grid from the distractors because my referring expression was too specific and focused on a single row. What I should have done is to look for a unique feature that is present in the target grid but not in the distractors. In this case, I should have noticed that the target grid has a unique "X" in the top{-}left corner, which is not present in the distractors. However, I also failed to notice that the target grid has a unique combination of "X"s in the second row from the top, which is not present in the distractors. Next time, I will look for such unique features and incorporate them into my referring expression. I will also make sure to consider multiple features and their combinations to increase the chances of distinguishing the target grid from the distractors. Additionally, I should have noticed that the target grid has a unique "X" in the bottom{-}right corner, which is not present in the distractors. I will make sure to include this feature in my referring expression as well.} \\
            }
        }
    }
    & & \\ \\

    \theutterance \stepcounter{utterance}  
    & & & \multicolumn{4}{p{0.6\linewidth}}{
        \cellcolor[rgb]{0.95,0.95,0.95}{
            \makecell[{{p{\linewidth}}}]{
                \texttt{\tiny{[P1$\rightarrow$P2]}}
                \texttt{Post{-}episode Reflexion (not sent to GM):} \\
\texttt{New plan: I failed to match the grid with the given expression because I didn't carefully examine the first row and the second row from the top of each grid. The expression states that the grid has an "X" in the top{-}left corner and a unique combination of "X"s in the second row from the top, where the second row from the top has an "X" in the second column and an "X" in the fifth column. I will re{-}examine the grids and select the one that meets this condition. Specifically, I will check the first row for exactly two "X"s and the second row from the top for an "X" in the second column and an "X" in the fifth column.} \\
            }
        }
    }
     \\ \\

\end{supertabular}
}

\end{document}
